% ============================================================================
% Chapter 4: PokeAgent Baseline Harness (Headers Only)
% ============================================================================

\newpage
\section{PokeAgent Baseline Harness}\label{sec:pokeagent}

\subsection{Early Development}\label{sec:early_development}

\subsubsection{Competition Era}\label{sec:competition_era}

\textit{[To be written. Early insights; simple starting point with four modules, no planning or advanced mapping/pathfinding. Figure of initial architecture.]}

\subsubsection{Towards Subagents}\label{sec:towards_subagents}

\textit{[To be written. Model capability increases (Gemini 2.5 $\to$ Gemini 3); CLI harnesses evolving with planning, tool use, and subagent abstractions. Figure of intermediate baseline from midterm.]}

\subsection{Early Performance}\label{sec:early_performance}

\subsubsection{Competition Era Results}\label{sec:competition_results}

\textit{[To be written. Competitors vs.\ PokeAgent; weakest baseline vs.\ Evelord (objective system + Porymap) vs.\ hyper-specific methods with limited generalization. Results figure from benchmark paper.]}

\subsubsection{Benchmark Paper Results}\label{sec:benchmark_results}

\textit{[To be written. Speedrunning harness vs.\ model comparisons; Gemini outperforms in visual reasoning; promise of CLI agents (Claude Code, Gemini CLI, Codex CLI); token burn tradeoffs. Results figure.]}

\subsubsection{Strongest Baseline}\label{sec:strongest_baseline}

\textit{[To be written. Fully fleshed-out system design developed for AutoEvolve experiments; orchestrator + subagent paradigm for context management; abstractions for coherence and localization. System design figure.]}

\subsection{System Design}\label{sec:system_design}

\subsubsection{Trajectory Window}\label{sec:trajectory_window}

\textit{[To be written. RLM-inspired session log; custom agents process context streams and pass analysis to orchestrator; context management and coherence maintenance.]}

\subsubsection{Memory}\label{sec:memory}

\textit{[To be written. Short-term (tool call results in 20-step window + store overviews) vs.\ long-term (CRUD-accessible at higher latency). Cognitive analogy.]}

\subsubsection{Skills}\label{sec:skills}

\textit{[To be written. Behavioral skills (textual strategies) and executable skills (Python in sandbox); skill lifecycle.]}

\subsubsection{Subagents}\label{sec:subagents}

\textit{[To be written. One-step subagents (reflect, verify, summarize) and delegated-loop subagents (battler, planner); tool whitelists and termination conditions.]}

\subsubsection{Sandbox and Code Execution}\label{sec:sandbox}

\textit{[To be written. Agent's ability to run\_code; sandboxed execution environment.]}
